\section*{الفصل \ref{ch:g11:redox} --- الأكسدة والاختزال}

\begin{solution}{exo:g11:redox:1}
\omterm{def:g11:redox:oxidant}{المؤكسد}: \omterm{def:g6:pure-substances-and-mixtures:species}{نوع كيميائي} قادر على اكتساب \omterm{def:g9:inside-the-atom:nucleus}{إلكترونات}. \omterm{def:g11:redox:oxidant}{المختزل}: \omterm{def:g6:pure-substances-and-mixtures:species}{نوع كيميائي} قادر على
التخلي عن \omterm{def:g9:inside-the-atom:nucleus}{إلكترونات}. \omterm{def:g11:redox:oxidation}{الأكسدة}: فقد \omterm{def:g9:inside-the-atom:nucleus}{إلكترونات}. \omterm{def:g11:redox:oxidation}{الاختزال}: اكتساب
\omterm{def:g9:inside-the-atom:nucleus}{إلكترونات}.
\end{solution}

\begin{solution}{exo:g11:redox:2}
$\ce{Fe^{2+}}/\ce{Fe}$؛ $\ce{Ag+}/\ce{Ag}$؛ $\ce{I2}/\ce{I-}$؛
$\ce{Fe^{3+}}/\ce{Fe^{2+}}$.
\end{solution}

\begin{solution}{exo:g11:redox:3}
\ce{Ag+ + e- <=> Ag}؛ \ce{Al^{3+} + 3e- <=> Al}؛
\ce{Fe^{3+} + e- <=> Fe^{2+}}؛ \ce{Cl2 + 2e- <=> 2Cl-}.
\end{solution}

\begin{solution}{exo:g11:redox:4}
\omterm{def:g11:redox:oxidation}{أكسدة} (تُفقد \omterm{def:g9:inside-the-atom:nucleus}{إلكترونات})؛ ويسلك الحديد سلوك \omterm{def:g11:redox:oxidant}{مختزل}.
\end{solution}

\begin{solution}{exo:g11:redox:5}
الخارصين يتأكسد وهو \omterm{def:g11:redox:oxidant}{المختزل}؛ و\ce{Cu^{2+}} يُختزل وهو
\omterm{def:g11:redox:oxidant}{المؤكسد}. إلكترونان لكل \omterm{def:g7:atoms-and-molecules:atom}{ذرة} خارصين.
\end{solution}

\begin{solution}{exo:g11:redox:6}
ننطلق من \ce{H2O2 -> H2O}: % ce-unbalanced-ok (step 1 of the method)
الأكسجين: ذرتا O على اليسار، إذن \ce{2H2O} على اليمين؛ الهيدروجين: 2 على
اليسار، و4 على اليمين، إذن \ce{2H+} على اليسار؛ الشحنة: $+2$ على
اليسار، و0 على اليمين، إذن إلكترونان على اليسار:
\ce{H2O2 + 2H+ + 2e- <=> 2H2O}.
\end{solution}

\begin{solution}{exo:g11:redox:7}
\ce{I2 + 2e- -> 2I-} و\ce{2S2O3^{2-} -> S4O6^{2-} + 2e-}؛ إلكترونان
في كل منهما، إذن \ce{I2 + 2S2O3^{2-} -> 2I- + S4O6^{2-}}.
\end{solution}

\begin{solution}{exo:g11:redox:8}
\ce{Cu -> Cu^{2+} + 2e-} و\ce{Ag+ + e- -> Ag} مضروبةً في 2:
\ce{Cu + 2Ag+ -> Cu^{2+} + 2Ag}. \omterm{def:g9:ions:ion}{وأيونات} \ce{Cu^{2+}} المتكوّنة زرقاء.
\end{solution}

\begin{solution}{exo:g11:redox:9}
\ce{C6H8O6 -> C6H6O6 + 2H+ + 2e-} و\ce{I2 + 2e- -> 2I-}:
\ce{C6H8O6 + I2 -> C6H6O6 + 2I- + 2H+}. وفيتامين C هو \omterm{def:g11:redox:oxidant}{المختزل}.
\end{solution}

\begin{solution}{exo:g11:redox:10}
\ce{Al -> Al^{3+} + 3e-} تعطي 3 \omterm{def:g9:inside-the-atom:nucleus}{إلكترونات}، و\ce{Cu^{2+} + 2e- -> Cu}
تأخذ 2؛ وأصغر عدد مشترك 6، فتُضرب الأولى في
2 والثانية في 3: \ce{2Al + 3Cu^{2+} -> 2Al^{3+} + 3Cu}.
\end{solution}

\begin{solution}{exo:g11:redox:11}
\ce{H2O2 + 2H+ + 2e- -> 2H2O} و\ce{Fe^{2+} -> Fe^{3+} + e-} مضروبةً في 2:
\ce{H2O2 + 2H+ + 2Fe^{2+} -> 2H2O + 2Fe^{3+}}.
\end{solution}

\begin{solution}{exo:g11:redox:12}
$n(\ce{Zn}) = 1.30/65.4 = \qty{1.99e-2}{mol}$. تترسب \omterm{def:g7:atoms-and-molecules:atom}{ذرة} نحاس واحدة
مقابل كل \omterm{def:g7:atoms-and-molecules:atom}{ذرة} خارصين تتأكسد، إذن $n(\ce{Cu}) = \qty{1.99e-2}{mol}$
و$m(\ce{Cu}) = 1.99\times 10^{-2} \times 63.5 = \qty{1.26}{g}$.
\end{solution}

\begin{solution}{exo:g11:redox:13}
\ce{Cr2O7^{2-} + 14H+ + 6e- -> 2Cr^{3+} + 7H2O}
و\ce{Fe^{2+} -> Fe^{3+} + e-} مضروبةً في 6:
\ce{Cr2O7^{2-} + 14H+ + 6Fe^{2+} -> 2Cr^{3+} + 6Fe^{3+} + 7H2O}.
ستة \omterm{def:g9:ions:ion}{أيونات} \ce{Fe^{2+}} لكل \omterm{def:g9:ions:ion}{أيون} ثنائي كرومات: $6 \times 0.010 = \qty{0.060}{mol}$.
\end{solution}

\begin{solution}{exo:g11:redox:14}
\ce{2ClO- + 4H+ + 2e- -> Cl2 + 2H2O} و\ce{2Cl- -> Cl2 + 2e-}؛ وبالجمع
والقسمة على 2: \ce{ClO- + Cl- + 2H+ -> Cl2 + H2O}. فالحمض يقدّم
\omterm{def:g9:ions:ion}{أيونات} \ce{H+}، والتفاعل يطلق ثنائي الكلور، وهو غاز سام:
ومن هنا التحذير من \omterm{def:g2:mixing-and-dissolving:mix}{خلط} المنتجين أبدًا.
\end{solution}

\begin{solution}{exo:g11:redox:15}
لا تستطيع \omterm{def:g9:inside-the-atom:nucleus}{الإلكترونات} أن تنتقل حرة عبر \omterm{def:g2:mixing-and-dissolving:dissolve}{المحلول}
(\cref{prop:g11:redox:no-free-electrons}): فالتي تفقدها \omterm{def:g7:atoms-and-molecules:atom}{ذرات} الخارصين لا
يأخذها إلا \omterm{def:g9:ions:ion}{أيون} \ce{Cu^{2+}} يلامس \omterm{def:g9:periodic-table-first-look:metal}{الفلز}. لذلك تتكوّن \omterm{def:g7:atoms-and-molecules:atom}{ذرة} النحاس
عند سطح الصفيحة، وتبقى هناك.
\end{solution}

\begin{solution}{pb:g11:redox:1}
\textbf{1.} $\ce{Cr2O7^{2-}}/\ce{Cr^{3+}}$ و$\ce{CH3COOH}/\ce{CH3CH2OH}$.

\textbf{2.} يتأكسد الإيثانول: فهو \omterm{def:g11:redox:oxidant}{المختزل}.

\textbf{3.} \ce{Cr2O7^{2-} + 14H+ + 6e- <=> 2Cr^{3+} + 7H2O}.

\textbf{4.} الكربون موزون (ذرتان في كل طرف). الأكسجين: ذرتان على
اليسار، وواحدة على اليمين، إذن \ce{H2O} على اليمين. الهيدروجين: أربع على
اليسار، و$6 + 2 = 8$ على اليمين، إذن \ce{4H+} على اليسار. الشحنة: $+4$ على
اليسار، و0 على اليمين، إذن أربعة \omterm{def:g9:inside-the-atom:nucleus}{إلكترونات} على اليسار:
\ce{CH3COOH + 4H+ + 4e- <=> CH3CH2OH + H2O}.

\textbf{5.} اثنا عشر إلكترونًا: نضرب الأولى في 2، والثانية معكوسةً
في 3:
\ce{2Cr2O7^{2-} + 3CH3CH2OH + 16H+ -> 4Cr^{3+} + 3CH3COOH + 11H2O}.

\textbf{6.} تُختزل \omterm{def:g9:ions:ion}{أيونات} \ce{Cr2O7^{2-}} البرتقالية إلى \omterm{def:g9:ions:ion}{أيونات} \ce{Cr^{3+}}
خضراء حيثما يصل إليها الإيثانول.

\textbf{7.} $M = 2\times 39.1 + 2\times 52.0 + 7\times 16.0 =
\qty{294.2}{g/mol}$.

\textbf{8.} $n = 2.0\times 10^{-3}/294.2 = \qty{6.8e-6}{mol}$.

\textbf{9.} ثلاثة \omterm{def:g7:atoms-and-molecules:molecule}{جزيئات} إيثانول لكل أيوني ثنائي كرومات:
$n = \tfrac{3}{2}\times 6.8\times 10^{-6} = \qty{1.0e-5}{mol}$.

\textbf{10.} $M(\ce{C2H6O}) = \qty{46.0}{g/mol}$، إذن
$m = 1.02\times 10^{-5}\times 46.0 = \qty{4.7e-4}{g}$، أي نحو
\qty{0.47}{mg}.

\textbf{11.} $0.25/0.47 \approx 0.53$: نحو \qty{53}{\percent} من
ثنائي الكرومات.

\textbf{12.} $0.53 \times 10 \approx \qty{5.3}{cm}$.

\textbf{13.} $\qty{0.25}{mg} \times 2100 = \qty{525}{mg}$ في لتر
الدم، أي نحو \qty{0.53}{g/L}.

\textbf{14.} يلقى النَّفَس الحبيبات عند المدخل أولًا؛ وثنائي الكرومات هناك
بفائض فيزيل الإيثانول كله، فيكتمل التفاعل في الجزء الأول من
الأنبوب، الذي يخضرّ تمامًا، قبل أن يصل أي إيثانول إلى
أبعد منه.

\textbf{15.} \omterm{def:g10:reaction-progress-table:limiting}{التفاعل تام}: \omterm{def:g9:ions:ion}{فأيونات} \ce{Cr^{3+}} الخضراء لا تعود
إلى ثنائي كرومات، فلا يستطيع الأنبوب المستعمل القياس من جديد.

% ledger: ghs:K2Cr2O7 (H350 among its hazard statements)
\textbf{16.} إنه شديد السمية (GHS06) وخطر جسيم على الصحة
(GHS08: من بيانات خطره «قد يسبب السرطان»)، فضلًا عن أنه
أكّال وسام للأحياء المائية؛ فلا يجوز أن يحويه جهاز
يستعمله الجمهور.

\textbf{17.} \ce{O2 + 4H+ + 4e- -> 2H2O}
و\ce{CH3CH2OH + H2O -> CH3COOH + 4H+ + 4e-}:
\ce{CH3CH2OH + O2 -> CH3COOH + H2O}.

\textbf{18.} أربعة \omterm{def:g9:inside-the-atom:nucleus}{إلكترونات}.

\textbf{19.} $n = 0.25\times 10^{-3}/46.0 = \qty{5.4e-6}{mol}$ من
الإيثانول؛ \omterm{def:g9:inside-the-atom:nucleus}{والإلكترونات} $4n = \qty{2.2e-5}{mol}$، أي
$2.2\times 10^{-5}\times 6.02\times 10^{23} = \num{1.3e19}$ \omterm{def:g9:inside-the-atom:nucleus}{إلكترون}.

\textbf{20.} $Q = 1.3\times 10^{19}\times 1.60\times 10^{-19} \approx
\qty{2.1}{C}$. فكل \omterm{def:g7:atoms-and-molecules:molecule}{جزيء} إيثانول يرسل أربعة \omterm{def:g9:inside-the-atom:nucleus}{إلكترونات} بالضبط
عبر السلك، فتكون الشحنة متناسبة مع كمية
الإيثانول: وقياسها قياس للإيثانول.
\end{solution}
